\documentclass[11pt]{article}

\usepackage[a4paper,margin=1in]{geometry}
\usepackage[T1]{fontenc}
\usepackage[utf8]{inputenc}
\usepackage{lmodern}
\usepackage{microtype}
\usepackage{amsmath,amssymb,amsthm,mathtools}
\usepackage{bm}
\usepackage{booktabs}
\usepackage{array}
\usepackage{tabularx}
\usepackage{enumitem}
\usepackage{xcolor}
\usepackage{hyperref}
\usepackage{listings}
\usepackage{caption}
\usepackage{fancyhdr}
\setlength{\headheight}{14pt}

\hypersetup{
  colorlinks=true,
  linkcolor=blue!50!black,
  citecolor=blue!50!black,
  urlcolor=blue!50!black,
  pdftitle={A Locally Separable 6-Parameter LQD Smile Model in Normalized Log-Returns},
  pdfauthor={OpenAI}
}

\definecolor{codebg}{RGB}{248,248,248}
\definecolor{codeframe}{RGB}{220,220,220}
\definecolor{codekw}{RGB}{0,76,153}
\definecolor{codecom}{RGB}{0,128,0}
\definecolor{codestr}{RGB}{153,0,0}

\lstdefinestyle{pycode}{
  language=Python,
  backgroundcolor=\color{codebg},
  frame=single,
  rulecolor=\color{codeframe},
  basicstyle=\ttfamily\footnotesize,
  keywordstyle=\color{codekw}\bfseries,
  commentstyle=\color{codecom},
  stringstyle=\color{codestr},
  showstringspaces=false,
  breaklines=true,
  columns=fullflexible,
  tabsize=4,
  xleftmargin=0.5em,
  xrightmargin=0.5em
}

\newcommand{\R}{\mathbb{R}}
\newcommand{\E}{\mathbb{E}}
\newcommand{\dd}{\mathrm{d}}
\newcommand{\Sig}{\Sigma}
\newcommand{\iv}{\sigma_{\mathrm{BS}}}
\newcommand{\vect}[1]{\bm{#1}}

\pagestyle{fancy}
\fancyhf{}
\lhead{Locally separable LQD smile model}
\rhead{\thepage}
\renewcommand{\headrulewidth}{0.3pt}

\title{A Locally Separable 6-Parameter LQD Smile Model \\ in Normalized Log-Returns}
\author{}
\date{March 2026}

\begin{document}
\maketitle
\thispagestyle{fancy}

\begin{abstract}
This note develops a six-parameter \emph{log quantile density} (LQD) smile model for a normalized log-return variable
\[
Z = \frac{X}{\alpha}, \qquad X = \log(S_T/F_T), \qquad \alpha = \sigma_{\mathrm{ref}}\sqrt{T},
\]
with three design goals: (i) the tails should be consistent with Lee/SVI-type linear variance wings, (ii) the six coefficients should each have a clear qualitative effect on the smile, and (iii) after a local rotation, the parameters should be locally separable on the actual implied-volatility curve $z \mapsto \iv(\alpha z,T)$. The construction is based on a fixed logarithmic tail skeleton, a six-function Fourier basis for the bounded LQD remainder, and a two-stage local rotation: first to trader-style smile observables, then to a weighted QR orthogonalization on the full $IV(z)$ curve. A NumPy-only reference implementation is included in the appendix and as a standalone Python file.
\end{abstract}

\tableofcontents

\section{Objective and notation}
The aim is to model a normalized slice with coordinate
\[
z = \frac{k}{\alpha}, \qquad k = \log(K/F_T), \qquad \alpha > 0,
\]
so that the \emph{shape} of the smile is expressed in comparable units across underlyings and maturities. We work with the normalized terminal return variable $Z$, its quantile function
\[
Q(u) := F_Z^{-1}(u), \qquad u \in (0,1),
\]
and its quantile density and log-quantile density
\[
q(u) := Q'(u), \qquad \ell(u) := \log q(u).
\]
By construction, $q(u) > 0$ for all $u \in (0,1)$ whenever $\ell$ is finite, so the quantile map is automatically increasing. This is one of the main practical attractions of the LQD representation: positivity of the quantile density is built in through the exponential map. Quantile-density methods and their estimation theory are classical; see Jones~\cite{Jones1992}.

Throughout the note, the six raw coefficients are collected in
\[
\theta = (p,c,m,s,\kappa,h)^\top \in \R^6,
\]
where the intended qualitative meanings are:
\begin{center}
\begin{tabular}{>{\raggedright\arraybackslash}p{2.0cm}p{10.2cm}}
\toprule
Coefficient & Intended raw interpretation \\
\midrule
$p$ & Put wing: left-tail constant in the LQD, hence far-left smile wing \\
$c$ & Call wing: right-tail constant in the LQD, hence far-right smile wing \\
$m$ & Belly level: mainly lifts or lowers the smile around the center \\
$s$ & Skew: first-order antisymmetric tilt around the center \\
$\kappa$ & Kurtosis / ATM butterfly: center curvature \\
$h$ & Shoulder: how fast the smile leaves the belly and connects to the wings \\
\bottomrule
\end{tabular}
\end{center}
Those meanings are exact only after a \emph{local} rotation described later; before rotation they are best understood as carefully designed, but still approximate, directions in LQD space.

\section{Why the boundary terms must be logarithmic}
\subsection{Lee/SVI tails imply logarithmic LQD singularities}
The boundary choice is dictated by the tail class. If the implied total variance has linear wings in the sense of Lee's moment formula~\cite{Lee2004}, then the log-return tails are of exponential type, consistent with the familiar SVI/SSVI asymptotics; see also Gatheral and Jacquier~\cite{GatheralJacquier2014}. Write the normalized tails as
\[
\mathbb{P}(Z > z) \asymp e^{-\Lambda_R z} \quad (z \to +\infty),
\qquad
\mathbb{P}(Z < z) \asymp e^{\Lambda_L z} \quad (z \to -\infty),
\]
with tail rates $\Lambda_L,\Lambda_R > 0$. Then the quantile asymptotics are
\[
Q(u) \sim \frac{1}{\Lambda_L}\log u \quad (u \downarrow 0),
\qquad
Q(u) \sim -\frac{1}{\Lambda_R}\log(1-u) \quad (u \uparrow 1),
\]
and differentiation yields
\[
q(u) \sim \frac{1}{\Lambda_L u} \quad (u \downarrow 0),
\qquad
q(u) \sim \frac{1}{\Lambda_R (1-u)} \quad (u \uparrow 1).
\]
Therefore the LQD must satisfy
\begin{equation}
\ell(u)
\sim -\log u - \log \Lambda_L \quad (u \downarrow 0),
\qquad
\ell(u)
\sim -\log(1-u) - \log \Lambda_R \quad (u \uparrow 1).
\label{eq:lqd-endpoints}
\end{equation}
The crucial consequence is that the coefficients of $-\log u$ and $-\log(1-u)$ are \emph{not} free. They must both be equal to $1$ if one wants to remain in the exponential-tail / linear-wing class. The free tail information is not in the singular coefficients, but in the \emph{endpoint constants}.

\subsection{Fixed tail skeleton}
For this reason the right starting point is
\begin{equation}
\ell(u) = -\log u - \log(1-u) + r(u),
\label{eq:lqd-skeleton}
\end{equation}
where $r(u)$ stays bounded at both endpoints. The bounded remainder controls the left and right constants and the interior shape, while the singular part enforces the correct tail class.

\section{A six-function Fourier basis for the bounded remainder}
\subsection{Raw model}
We write the bounded remainder as a six-term Fourier expansion,
\begin{equation}
r(u;\theta) = \sum_{j=1}^6 \theta_j \, \Phi_j(u),
\qquad
\ell(u;\theta) = -\log u - \log(1-u) + r(u;\theta),
\label{eq:raw-lqd-model}
\end{equation}
with the six basis functions chosen so that two of them control the endpoint constants and the other four isolate the first few center jets.

\subsection{Wing modes}
Define
\begin{align}
\Phi_P(u)
&= \frac{3}{16} + \frac{3}{8}\cos(\pi u) + \frac{1}{4}\cos(2\pi u)
   + \frac{1}{8}\cos(3\pi u) + \frac{1}{16}\cos(4\pi u),
\label{eq:PhiP}\\
\Phi_C(u)
&= \Phi_P(1-u)
 = \frac{3}{16} - \frac{3}{8}\cos(\pi u) + \frac{1}{4}\cos(2\pi u)
   - \frac{1}{8}\cos(3\pi u) + \frac{1}{16}\cos(4\pi u).
\label{eq:PhiC}
\end{align}
These satisfy
\[
\Phi_P(0)=1,\ \Phi_P(1)=0,
\qquad
\Phi_C(0)=0,\ \Phi_C(1)=1,
\]
while around the center $u = \tfrac12 + t$ they start only at cubic order,
\[
\Phi_P\Bigl(\tfrac12+t\Bigr) = -\frac{\pi^3}{2}t^3 + O(t^4),
\qquad
\Phi_C\Bigl(\tfrac12+t\Bigr) = +\frac{\pi^3}{2}t^3 + O(t^4).
\]
So they mainly change the endpoint constants and hence the far wings, while leaving the ATM level, ATM slope, and ATM curvature almost untouched at leading order.

With coefficients $(p,c)$ in front of $(\Phi_P,\Phi_C)$, the endpoint expansions become
\[
\ell(u) = -\log u + p + O(u) \quad (u \downarrow 0),
\qquad
\ell(u) = -\log(1-u) + c + O(1-u) \quad (u \uparrow 1),
\]
so the normalized tail rates are
\begin{equation}
\Lambda_L = e^{-p},
\qquad
\Lambda_R = e^{-c}.
\label{eq:tail-rates}
\end{equation}
This is the exact sense in which $p$ and $c$ are the put-wing and call-wing parameters.

\subsection{Belly modes}
For the interior modes we use sine combinations that vanish at both endpoints and isolate the first few center jets:
\begin{align}
\Phi_0(u)
&= \frac{75}{64}\sin(\pi u) + \frac{25}{128}\sin(3\pi u) + \frac{3}{128}\sin(5\pi u), \\
\Phi_1(u)
&= -\frac{8\sin(2\pi u) + \sin(4\pi u)}{12\pi}, \\
\Phi_2(u)
&= \frac{34\sin(\pi u) + 39\sin(3\pi u) + 5\sin(5\pi u)}{96\pi^2}, \\
\Phi_4(u)
&= \frac{2\sin(\pi u) + 3\sin(3\pi u) + \sin(5\pi u)}{16\pi^4}.
\end{align}
Near the center $u = \tfrac12 + t$ they satisfy
\begin{equation}
\Phi_0 = 1 + O(t^6),
\qquad
\Phi_1 = t + O(t^5),
\qquad
\Phi_2 = t^2 + O(t^6),
\qquad
\Phi_4 = t^4 + O(t^6).
\label{eq:center-jets}
\end{equation}
This is the key rationale for the choice: the four interior coefficients align, as much as possible in a low-dimensional model, with level, skew, curvature, and shoulder-width directions.

\subsection{Summary table}
Collecting all six modes,
\[
\Phi(u) = \bigl(\Phi_P(u),\Phi_C(u),\Phi_0(u),\Phi_1(u),\Phi_2(u),\Phi_4(u)\bigr),
\qquad
\theta = (p,c,m,s,\kappa,h)^\top,
\]
so that
\[
\ell(u;\theta) = -\log u - \log(1-u) + \Phi(u)\theta.
\]
The qualitative role of each mode is summarized in Table~\ref{tab:raw-basis-summary}.

\begin{table}[h]
\centering
\caption{Raw six-function LQD basis and its intended role.}
\label{tab:raw-basis-summary}
\begin{tabularx}{\textwidth}{>{\raggedright\arraybackslash}p{2.1cm} >{\raggedright\arraybackslash}p{4.6cm} X}
\toprule
Mode & Main effect & Why it is separated \\
\midrule
$\Phi_P$ & Put wing / left tail constant & $\Phi_P(0)=1$, $\Phi_P(1)=0$, and its center jet starts only at cubic order \\
$\Phi_C$ & Call wing / right tail constant & $\Phi_C(1)=1$, $\Phi_C(0)=0$, and its center jet starts only at cubic order \\
$\Phi_0$ & Belly level & Equal to $1+O(t^6)$ at the center and zero at both endpoints \\
$\Phi_1$ & Skew & First-order antisymmetric center jet, zero at both endpoints \\
$\Phi_2$ & ATM curvature / butterfly & Second-order symmetric center jet, zero at both endpoints \\
$\Phi_4$ & Shoulder / wing-connection & Fourth-order symmetric center jet, zero at both endpoints \\
\bottomrule
\end{tabularx}
\end{table}

\section{Stable numerical construction in logistic coordinates}
A practical advantage of the fixed tail skeleton~\eqref{eq:lqd-skeleton} is that it disappears exactly under the logistic change of variables
\[
x = \log\frac{u}{1-u},
\qquad
u(x) = \frac{1}{1+e^{-x}},
\qquad
\frac{\dd u}{\dd x} = u(1-u).
\]
Since
\[
q(u) = e^{\ell(u)} = \frac{e^{r(u)}}{u(1-u)},
\]
we obtain the remarkably simple identity
\begin{equation}
\frac{\dd Q}{\dd x}
= q(u(x))\frac{\dd u}{\dd x}
= e^{r(u(x))}.
\label{eq:dQdx}
\end{equation}
Thus the endpoint singularities are regularized \emph{exactly}. Numerically, this is the right coordinate in which to integrate the quantile. Once $Q$ is constructed, it is convenient to center it by imposing
\[
Q\Bigl(\tfrac12\Bigr)=0,
\]
which corresponds to setting the median of $Z$ to zero. Any remaining location shift is then handled later through the risk-neutral martingale constraint.

\subsection{Risk-neutral embedding and smile extraction}
Let
\[
X = \mu + \alpha Z,
\]
where the positive scale $\alpha$ is the normalization used on the strike axis, typically a reference ATM total standard deviation. The martingale condition $\E[e^X]=1$ determines
\begin{equation}
\mu(\theta)
= -\log \E\left[e^{\alpha Z}\right]
= -\log\int_0^1 e^{\alpha Q(u;\theta)}\,\dd u.
\label{eq:martingale-shift}
\end{equation}
Then the forward call price at log-moneyness $k$ is
\begin{equation}
C(k;\theta)
= \int_0^1 \Bigl(e^{\mu(\theta)+\alpha Q(u;\theta)} - e^k\Bigr)^+\,\dd u.
\label{eq:call-price}
\end{equation}
Finally, the actual smile on the normalized abscissa is
\begin{equation}
\Sig(z;\theta) := \iv(\alpha z, T; \theta),
\label{eq:smile-map}
\end{equation}
namely the Black-Scholes implied volatility extracted from the price at $k = \alpha z$ and maturity $T$.

\section{Why a local rotation is needed}
The raw parameters $(p,c,m,s,\kappa,h)$ are well designed, but the map
\[
\theta \longmapsto \Sig(\cdot;\theta)
\]
is nonlinear and couples the raw directions. In other words, a unit increase in $m$ does not move only the belly level once prices are translated into implied volatilities; similarly, a pure raw skew move slightly perturbs the wings and the minimum. Therefore, if the objective is \emph{local separability on the actual IV curve}, one should rotate the raw coefficients around a reference slice.

There are two good reasons to perform the rotation in two stages rather than one:
\begin{enumerate}[leftmargin=2.2em]
\item first align the raw coefficients with \emph{interpretable smile observables};
\item then orthogonalize those observable directions on the full curve with a weighted QR.
\end{enumerate}
A direct QR of the raw basis would produce orthogonal directions, but it would discard the trader intuition attached to the six parameters.

\section{Stage 1: local trader coordinates}
Fix a reference slice $\theta_*$ and a normalized strike grid
\[
z_1,\dots,z_N.
\]
Define the smile vector
\[
\sigma(\theta)
:= \bigl(\Sig(z_1;\theta),\dots,\Sig(z_N;\theta)\bigr)^\top \in \R^N.
\]
Its Jacobian at the reference point is
\begin{equation}
J := \frac{\partial \sigma}{\partial \theta}(\theta_*) \in \R^{N\times 6}.
\label{eq:J-def}
\end{equation}

\subsection{Six local smile observables}
We now choose six linearized smile observables that match the intended meaning of the parameters:
\begin{align}
y_1 &= \Sig(z_{\min}^*), \\
y_2 &= \Sig_z(0), \\
y_3 &= \Sig_{zz}(0), \\
y_4 &= S_P := \frac{\Sig(-z_2)-\Sig(-z_1)}{z_1-z_2}, \\
y_5 &= S_C := \frac{\Sig(z_2)-\Sig(z_1)}{z_2-z_1}, \\
y_6 &= Sh := \frac{\Sig(z_s)+\Sig(-z_s)}{2} - \Sig(0) - \frac{z_s^2}{2}\Sig_{zz}(0).
\label{eq:observable-defs}
\end{align}
Here $z_{\min}^*$ is the minimum location of the \emph{reference} smile, frozen once the linearization point is chosen. This freezing is important: if the minimum location were re-optimized under perturbations, the observable would become nonlinear and would re-entangle skew and curvature.

The interpretation of the six observables is:
\[
(y_1,y_2,y_3,y_4,y_5,y_6)
=
(\text{min-level},\ \text{ATM skew},\ \text{ATM curvature},\ \text{put wing},\ \text{call wing},\ \text{shoulder}).
\]
The shoulder observable subtracts the quadratic center contribution so that it reacts mainly to the belly-to-wing transition rather than to plain ATM curvature.

In numerical work we realize these observables as linear combinations of smile values on the grid $z_1,\dots,z_N$. Hence there exists an operator
\[
O \in \R^{6\times N}
\]
such that
\[
y(\theta) = O\,\sigma(\theta),
\]
at least after freezing $z_{\min}^*$ and choosing the finite-difference stencils.
Then
\begin{equation}
K := \frac{\partial y}{\partial \theta}(\theta_*) = OJ \in \R^{6\times 6}.
\label{eq:K-def}
\end{equation}
The \emph{trader coordinates} are now defined by
\begin{equation}
c := K(\theta-\theta_*).
\label{eq:trader-coords}
\end{equation}
By construction,
\begin{equation}
\delta y = \delta c + O(\|\delta\theta\|^2).
\label{eq:delta-y-delta-c}
\end{equation}
So, to first order, $c_1$ is a pure minimum-level move, $c_2$ a pure ATM skew move, and so on.

\paragraph{Remark.}
Because the construction is on the actual IV curve, $y_1$ is a \emph{minimum-IV} observable. If one wants a closer analogue of the SVI-JW minimum-variance parameter, replace $y_1$ by the minimum of total variance or total standard deviation; the methodology is unchanged.

\section{Stage 2: weighted QR orthogonalization on the smile curve}
The trader coordinates are interpretable but not orthogonal on the whole curve. To obtain a numerically convenient basis for optimization and regularization, choose a diagonal positive weight matrix
\[
W = \operatorname{diag}(w_1,\dots,w_N),
\qquad w_i > 0,
\]
for example vega weights, bid-ask inverse-variance weights, or a center-heavy proxy. Consider
\begin{equation}
A := W^{1/2} J K^{-1} \in \R^{N\times 6}.
\label{eq:A-def}
\end{equation}
Perform a reduced QR decomposition,
\begin{equation}
A = QR,
\qquad Q^\top Q = I_6,
\qquad R \text{ upper triangular}.
\label{eq:qr-decomp}
\end{equation}
Then define the \emph{orthogonal curve coordinates}
\begin{equation}
b := Rc = RK(\theta-\theta_*).
\label{eq:b-def}
\end{equation}
Substituting~\eqref{eq:b-def} into the linearized smile map yields
\begin{equation}
\delta\sigma
\approx J\,\delta\theta
= W^{-1/2}Q\,\delta b.
\label{eq:orthogonal-smile-moves}
\end{equation}
So the six columns of $Q$ are orthonormal local smile moves in the weighted IV metric. This is the exact sense in which the $b$-coordinates are \emph{locally separable on the actual $IV(z)$ curve}.

\subsection{Three coordinate systems}
The construction produces three useful coordinate systems:
\begin{center}
\begin{tabular}{>{\raggedright\arraybackslash}p{2.3cm} >{\raggedright\arraybackslash}p{3.2cm} p{7.0cm}}
\toprule
System & Symbol & Role \\
\midrule
Raw LQD coefficients & $\theta$ & Natural parametric basis in quantile-density space \\
Trader coordinates & $c = K(\theta-\theta_*)$ & Locally equal to the chosen six IV observables \\
Orthogonal curve coordinates & $b = Rc$ & Locally orthogonal on the weighted smile curve \\
\bottomrule
\end{tabular}
\end{center}
In practice, the clean workflow is: \emph{report} $c$, \emph{optimize} in $b$, and keep $\theta$ only as the internal representation.

\section{Rotated LQD basis functions}
Since
\[
\theta-\theta_* = K^{-1}c = K^{-1}R^{-1}b,
\]
the LQD can be rewritten as
\begin{equation}
\ell(u;b)
= \ell(u;\theta_*) + \sum_{j=1}^6 b_j\,\Psi_j(u),
\label{eq:rotated-lqd}
\end{equation}
where the rotated basis functions are
\begin{equation}
\Psi_j(u) = \sum_{i=1}^6 \Phi_i(u)\,\bigl(K^{-1}R^{-1}\bigr)_{ij}.
\label{eq:rotated-basis}
\end{equation}
These $\Psi_j$ are still basis functions in LQD space, but they are now aligned with locally orthogonal smile moves. In other words, the raw Fourier basis gives the right \emph{functional scaffold}, while the local QR gives the right \emph{market geometry}.

\section{Recommended practical choices}
A robust default setup is:
\begin{itemize}[leftmargin=1.8em]
\item normalized strike grid $z \in [-3,3]$ with roughly $40$ to $60$ points;
\item center finite-difference step $\delta = 0.25$;
\item wing intervals $[z_1,z_2] = [1.5,2.5]$;
\item shoulder point $z_s = 1.0$;
\item weights proportional to option vega divided by squared bid-ask spread, if market data are available.
\end{itemize}
If one only wants a model-space orthogonalization, a smooth center-heavy choice such as $w_i = 1/(1+z_i^2)$ works well.

\section{Calibration workflow}
The recommended calibration loop for a single slice is:
\begin{enumerate}[leftmargin=1.9em]
\item Fit the raw six-parameter LQD model $\theta$ to the slice.
\item Freeze the resulting fit as $\theta_*$ and build $J$, $O$, $K$, and the QR factors.
\item Re-express perturbations in orthogonal coordinates $b$ for optimization or temporal smoothing.
\item Report the slice in trader coordinates $c$, which have the clearest interpretation.
\item Refresh the local rotation whenever the reference slice moves materially; the entire construction is local by design.
\end{enumerate}

\section{Why this construction is useful}
The methodology solves three problems at once:
\begin{enumerate}[leftmargin=1.9em]
\item \textbf{Tail correctness.} The fixed $-\log u$ and $-\log(1-u)$ singularities enforce the exponential tail class associated with Lee/SVI-type linear wings.
\item \textbf{Interpretability.} The six raw Fourier modes are already designed to act like put wing, call wing, belly level, skew, curvature, and shoulder.
\item \textbf{Local separability on the actual smile.} The two-stage local rotation makes the parameters act independently on the actual implied-volatility curve rather than merely on the LQD.
\end{enumerate}
In that sense, the model plays the same role for LQD-based smiles that SVI-JW plays for raw SVI: the representation separates level from shape as much as possible, and then the local coordinates are chosen to match directly interpretable smile moves.

\appendix
\section{NumPy reference implementation}
The accompanying Python file implements the full methodology:
\begin{itemize}[leftmargin=1.8em]
\item the six raw LQD basis functions,
\item stable quantile construction in logistic coordinates,
\item risk-neutral martingale centering,
\item option pricing and Black-Scholes inversion,
\item the observable operator $O$,
\item the Jacobian matrices $J$ and $K$,
\item the weighted QR rotation,
\item the rotated LQD basis functions.
\end{itemize}
For convenience the complete source is included below.

\lstset{style=pycode}
\lstinputlisting{lqd_local_qr.py}

\begin{thebibliography}{9}

\bibitem{Jones1992}
M. C. Jones,
\emph{Estimating densities, quantiles, quantile densities and density quantiles},
Annals of the Institute of Statistical Mathematics, 44(4):721--727, 1992.

\bibitem{Lee2004}
Roger W. Lee,
\emph{The moment formula for implied volatility at extreme strikes},
Mathematical Finance, 14(3):469--480, 2004.

\bibitem{GatheralJacquier2014}
Jim Gatheral and Antoine Jacquier,
\emph{Arbitrage-free SVI volatility surfaces},
Quantitative Finance, 14(1):59--71, 2014.

\end{thebibliography}

\end{document}
